\documentclass[10pt,a4paper]{amsart}
\usepackage[margin=19mm]{geometry}
\usepackage{amsmath,amssymb,mathtools}
\usepackage[scr=rsfs,cal=euler]{mathalfa}
\usepackage{xfrac}
\usepackage[unicode,hidelinks,bookmarks=false]{hyperref}
\newcommand{\ie}{i.e.\ }
\newcommand{\cf}{cf.\ }
\newcommand{\resp}{respectively\ }
\newcommand{\Subsection}{\subsection}
\providecommand{\nobreakdash}{\nobreak\hbox{-}}
\setlength{\emergencystretch}{2em}
\multlinegap0pt
\usepackage{amssymb}
\usepackage{mathtools}
\usepackage{xcolor}
\newtheorem{lemma}{Lemma}
\newtheorem{proposition}{Proposition}
\newcommand{\N}{{\mathbb N}}
\newcommand{\R}{{\mathbb R}}
\newcommand{\Z}{{\mathbb Z}}
\title{Asymptotic safety regions for Gabor frames generated by Hermite functions}
\author{Markus Faulhuber, Irina Shafkulovska, Ilya Zlotnikov}
\date{Source: arXiv:2609.01296v1 (CC BY 4.0)}
\begin{document}
\maketitle

\noindent\emph{Source and licence.} Asymptotic safety regions for Gabor frames generated by Hermite functions. Version arXiv:2609.01296v1 (CC BY 4.0), distributed by arXiv under the Creative Commons Attribution 4.0 International licence, http://creativecommons.org/licenses/by/4.0/, as recorded in the arXiv licence metadata for this exact version and shown on the abstract page https://arxiv.org/abs/2609.01296v1. Full licence notice: \texttt{licenses/arxiv-2609.01296-CC-BY-4.0.txt}.\par\smallskip

\noindent\textit{Editorial scope.} Counted unit: the article's proposition pro:Janssen\_main (source0.tex lines 709--718). The article's complete original proof of it is delivered with it (source0.tex lines 720--905, one \texttt{proof} environment of 186 lines). No mathematics is written, repaired or re-derived: the statement and the proof are the article's own lines, in the article's own notation, and no retained line is substituted or re-flowed.  Retained uncounted as the source-local context the proof needs: lines 345--345; lines 364--366; lines 393--399; lines 463--472; lines 469--471; lines 476--478; lines 527--529; lines 632--636.  Nothing was omitted from any retained span, so the package carries no omission record.  No retained line contains a citation, so the wrapper carries no bibliography and no bibliography entry.  No retained line contains a figure environment, an image-inclusion command or drawing code, so no image is delivered and the package declares no asset dependency.  Editorial formatting only, in addition: an AMS \texttt{amsart} preamble that restates the article's own declarations and macros used by the retained lines, namely the environments \texttt{lemma}, \texttt{proposition} on separate counters, together with the standard packages those lines need.  The retained environments print in this document as Lemma 1, Proposition 1 in order of appearance, whereas the article prints the same items with the numbering of its own full document.  The complete native source of the article as distributed in the arXiv e-print for this exact version is bundled unchanged as \texttt{sources/209094/source0.tex}.
\subsection{Upper bounds on Laguerre polynomials}\label{sec:upper_bounds_for_laguerre}

\begin{equation}\label{eq:q+s=nu}
    q + s = (\sqrt{n+1}-\sqrt{n})^2 + (\sqrt{n+1}+\sqrt{n})^2 = 4n+2 =: \nu.
\end{equation}

\begin{equation}\label{eq:estimate}
    |\mathcal{L}^{0}_n(x)|e^{-x/2}\leq \begin{cases}
        \sqrt{\frac{s-q}{(s-x)(x-q)}}, &\quad  q\leq x\leq s\\
        C n^{-1/3} &\quad  \nu -\big(\tfrac{4\nu}{3}\big)^{1/3} \leq x\leq \nu +\big(\tfrac{4\nu}{3}\big)^{1/3} \\
        C x^{-1/2} n^{1/6} e^{-R(s,x)/2}, &\quad  x\geq s.
    \end{cases}
\end{equation}

\begin{lemma}\label{lem:ellipse_lattice_count}
Let $0< a \leq b$ and $r\geq0$. Define an ellipse
    \begin{equation}
        E_r = \left \{(x,y)\in \R^2 :\frac{x^2}{b^2}+\frac{y^2}{a^2}\leq \frac{r}{\pi} \right \}.
    \end{equation}
There exists an absolute constant $C>0$ such that for every $0 \leq r_1 \leq r_2$ we have
\begin{equation}\label{eq:ellipse_difference}
     \#\left(( E_{r_2} \setminus E_{r_1})\cap \Z^2 \right) \leq ab(r_2-r_1) + C(a\sqrt{r_2} + b \sqrt{r_2-r_1} + 1).
\end{equation}
\end{lemma}

\begin{equation}
     \#\left(( E_{r_2} \setminus E_{r_1})\cap \Z^2 \right) \leq ab(r_2-r_1) + C(a\sqrt{r_2} + b \sqrt{r_2-r_1} + 1).
\end{equation}

\begin{equation}\label{eq:ellipse}
     \#\left( E_{r}\cap \Z^2 \right) \leq abr + C((a+b)\sqrt{r} + 1).
\end{equation}

    \begin{equation}\label{eq:diagonal_safety_region}
           S_n(\delta,\eta) := \left\{ (a,b) \in \R^2_+ : \min\{a,b\}\geq n^{-\frac12-\delta}, ab \leq n^{-\frac{2}{3}-\eta} \right\} 
        \end{equation}

    \begin{equation}\label{eq:X_uniform_bounds}
        \pi (k^2 + l^2) n^{\frac{1}{3}+2(\eta - \delta)} = \pi (k^2 + l^2) n^{1-\varepsilon}
        \leq X_{k,l}(a,b)
        \leq \pi (k^2 + l^2) n^{1+2\delta}.
    \end{equation}

\begin{proposition}\label{pro:Janssen_main}
    Let $0< \delta < \eta < \frac{1}{3}$, and $S_n(\delta,\eta)$ be defined by \eqref{eq:diagonal_safety_region}. Let $K$ denote the set
    \begin{equation}\label{eq:K_set_def}
        K := K(a,b):= \left\{(k,l)\in\Z^2\setminus\{(0,0)\}:\pi \left( \frac{k^2}{b^2}+\frac{l^2}{a^2} \right) \leq 10n\right\}.
    \end{equation}
    Then
    \begin{equation}
        \lim\limits_{n\to\infty} \sup\limits_{(a,b) \in S_n(\delta, \eta)} \sum_{(k,l)\in K} |\mathcal{L}^{0}_n(X_{k,l}(a,b))| e^{-\frac{X_{k,l}(a,b)}{2}} = 0.
    \end{equation}
\end{proposition}

\begin{proof}
    Without loss of generality, we may assume that $a \leq b$, $(a,b) \in S_n(\delta, \eta)$, such that 
    \begin{equation}\label{eq:X_uniform_bounds_reminder}
        \pi (k^2 + l^2) n^{\frac{1}{3}+2(\eta - \delta)} 
        \leq X_{k,l}(a,b)
        \leq \pi (k^2 + l^2) n^{1+2\delta}.
    \end{equation}
    In particular, this implies that all values $ X_{k,l}(a,b)$ with $(k,l) \in K$ satisfy
    \begin{equation}\label{eq:min_point}
        X_{k,l}(a,b) \geq \pi n^{\frac{1}{3}+2(\eta - \delta)}>\frac{1}{4n}\geq q.
    \end{equation}
    In what follows,  for $X_{k,l}(a,b)$ we will use only these bounds independent of $a$ and $b$.
    
    Next, we fix some parameters $p$ and $t$ satisfying conditions
    \begin{equation}\label{eq:p_t_cond}
        0<t<2/3, \,\, \tfrac{1}{3}<p<1,\quad \text{and} \quad     p+t > \frac{5}{3}-2(\eta - \delta).   
    \end{equation}
    Consider the partition $K = \bigcup_{j=1}^5 K_j$ defined in the following way: 
    \begin{align}
        K_1 &= \{(k,l)\in K: q\leq \pi(\tfrac{k^2}{b^2}+\tfrac{l^2}{a^2})\leq q+n^{t}\}; \\
        K_2 &= \{(k,l)\in K: q+n^{t}< \pi(\tfrac{k^2}{b^2}+\tfrac{l^2}{a^2})\leq s-n^{p}\}; \\
        K_3 &= \{(k,l)\in K: s-n^{p}< \pi(\tfrac{k^2}{b^2}+\tfrac{l^2}{a^2})\leq s-n^{1/3}\}; \\
        K_4 &= \{(k,l)\in K: s-n^{1/3} <\pi(\tfrac{k^2}{b^2}+\tfrac{l^2}{a^2})\leq s\}; \\
        K_5 &= \{(k,l)\in K: s <\pi(\tfrac{k^2}{b^2}+\tfrac{l^2}{a^2})\leq {10}n\}.
    \end{align} 
    Note that, since $0<t,p<1$ and $4n \leq s$, we have $q+n^t<s-n^p$ for all sufficiently large $n\in\N$. Moreover, $p>\frac13$ implies $s-n^p<s-n^{1/3}$ and $s \leq 4n+2 < 10n$, $n \in \N$. Thus, the above partition is well-defined (for $n$ sufficiently large) and the sets $K_j$ are pairwise disjoint.
        
    We now estimate the summands within each subset $K_j, \, j=1, \dots,5$ using the estimates for the Laguerre polynomials described in Subsection~\ref{sec:upper_bounds_for_laguerre}. The layer $K_3$ is more delicate than the others, so we first deal with
    \begin{equation}
        J_1:=\sum_{(k,l)\in K_1\cup K_2 \cup K_4\cup K_5} |\mathcal{L}_n^{0}(X_{k,l}(a,b))|e^{-\frac{X_{k,l}(a,b)}{2}}.
    \end{equation}
    
    For the layers $K_1$ and $K_2$, we use the first estimate in \eqref{eq:estimate} and ~\eqref{eq:X_uniform_bounds}.
    To apply the second estimate in \eqref{eq:estimate} for the layer $K_4$, we need to verify $\nu-(\tfrac{4\nu}{3})^{1/3}<s-n^{1/3}$. Recall from \eqref{eq:q+s=nu} that $\nu =s+q$ and, indeed,
      \begin{equation}
          \nu-(\tfrac{4\nu}{3})^{1/3} \leq s+q-(\tfrac{16n}{3})^{1/3} \leq s - n^{1/3} + \frac{1}{4n} +(1-(\tfrac{16}{3})^{1/3})n^{1/3}
          <s-n^{1/3}
        \end{equation}
    for all $n\geq 3$.
    
    The last envelope in \eqref{eq:estimate}, given by $C x^{-1/2} n^{1/6} e^{-R(s,x)/2}$, where $C$ is a constant independent of $x$ and $n$, is strictly monotonically decreasing in $x$ for $x \geq s$.
    Denoting the number of points in $K_j$ by $\# K_j$ we obtain the following estimate:
    \begin{equation}\label{eq:sum_K123}
        \begin{split}
            J_1
            & \leq \quad \#K_1\cdot \sqrt{\frac{s-q}{(s-(q+n^{t}))(\pi n^{\frac{1}{3}+2(\eta -\delta)}-q)}} \\
            & \quad + \#K_2 \cdot \sqrt{\frac{s-q}{(s-(s-n^{p}))((q+n^{t})-q)}}\\
            & \quad +C \left(\#K_4\cdot n^{-1/3} +\# K_5  \cdot s^{-1/2} n^{1/6} e^{-R(s,s)/2}\right) \\
            &= \#K_1 \cdot \sqrt{\frac{s-q}{(s-q-n^{t})(\pi n^{\frac{1}{3}+2(\eta -\delta)}-q)}} 
             + \#K_2 \cdot \sqrt{\frac{s-q}{n^{p+t}}} \\
            & \quad + C \left(\#K_4\cdot n^{-1/3} + \tfrac{1}{2} \cdot \# K_5 \cdot n^{-1/2+1/6}\right).
        \end{split}
    \end{equation}
    We proceed with simplifying this expression. It holds that $s-q = 4\sqrt{n(n+1)}\in (4n,4n+2)$, and, thus we have
    \begin{equation}
        n^{t}= 4n^{t}-3n^{t} \leq 3(4n-n^{t})
        \leq 3(s-q-n^{t}), 
        \quad 0 < t < 1,
    \end{equation}
    and consequently
    \begin{equation}
        \frac{s-q}{s-q-n^t} = 1+\frac{n^t}{s-q-n^t}\leq 4.
    \end{equation}
    It is clear that for sufficiently large $n$ we have $\pi n^{\frac{1}{3}+2(\eta -\delta)}-q>2n^{\frac{1}{3}+2(\eta -\delta)}$. Thus, we can now further estimate $J_1$ by
    \begin{equation}\label{eq:sum_update}
        \begin{split}
          J_1
            & \leq \#K_1 \cdot \sqrt{\tfrac{4}{2 n^{\frac{1}{3}+2(\eta -\delta)}}} 
            + 3\cdot \#K_2\cdot \sqrt{\tfrac{n}{n^{p+t}}} \\
            & \quad + C \left(\#K_4\cdot n^{-1/3} +\tfrac{1}{2}\# K_5 \cdot n^{-1/3}\right) \\
            &\leq 
            \sqrt{2} \cdot \#K_1 \cdot n^{-\frac{1}{6}-(\eta-\delta)} 
            + 3 \cdot \#K_2\cdot n^{(1-p-t)/2}\\
            & \quad + C \left(\#K_4\cdot n^{-1/3} +\tfrac{1}{2}\# K_5 \cdot n^{-1/3}\right).
        \end{split}
    \end{equation}
    To obtain an upper bound for $J_1$, it remains to estimate $\#K_j$, $j=1,2,4,5$. To this end, we use Lemma~\ref{lem:ellipse_lattice_count}. For $j=2$ and $j=5$ we use~\eqref{eq:ellipse} and the rough estimate  
    \begin{align}\label{eq:K25ab}
        \# K_j
         \leq \#K \leq \#(E_{10n}\cap \Z^2)
         \leq C_j (ab \, n + b \, n^{1/2} + 1),
    \end{align}
    where $C_j$, $j=2,5$, are universal constants, independent of $a, \, b, \, n$.
    Next, the same Lemma, combined with basic estimates, gives us
    \begin{equation}\label{eq:K1ab}
        \# K_1 \leq \#((E_{q+n^t} \setminus E_{q})\cap \Z^2) \leq C_1 ( ab n^{t} + a(q+n^t)^{1/2} + bn^{t/2}+1);
    \end{equation}
    \begin{equation}\label{eq:K4ab}
        \# K_4 \leq \#((E_{s} \setminus E_{s-n^{1/3}})\cap \Z^2)
        \leq C_4 ( ab n^{1/3} + a n^{1/2} + bn^{1/6}+1),
    \end{equation}
    with $C_1$ and $C_4$ again independent of $a, \, b, \, n$.

    Before returning to the estimate of $J_1$, we prefer to eliminate the dependence on $a$ and $b$ in the estimates of the cardinalities of the sets $K_j$. Recall that $0<\delta<\eta<1/3$, $0<t<1$, and 
    \begin{equation}\label{eq:ab_bounds_revisited}
        ab \leq n^{-2/3-\eta}, \quad a \leq \sqrt{ab} \leq n^{-1/3 - \eta/2}, \quad b = \frac{ab}{\min\{a,b\}} \leq \frac{n^{-2/3-\eta}}{n^{-1/2-\delta}} \leq n^{-1/6+\delta - \eta}.  
    \end{equation}
    Combining this with~\eqref{eq:K25ab},\eqref{eq:K1ab}, and \eqref{eq:K4ab}, we arrive at
     \begin{equation}\label{eq:K1245}
         \begin{split}
             \#K_j &\leq C_j(n^{1/3-\eta} +n^{1/3+\delta - \eta} +1)\\
             & \leq C_j n^{1/3+\delta-\eta}, \quad j=2,5;\\
             \#K_1 & \leq C_1(n^{t-2/3-\eta} + n^{t/2 - 1/3 - \eta/2}+n^{t/2-1/6+\delta-\eta}+1)\\
             & \leq C_1(n^{t/2-1/6+\delta-\eta}+1); \\
             \#K_4 & \leq C_4(n^{-1/3-\eta} + n^{1/6 - \eta/2}+n^{\delta-\eta}+1)\\
             & \leq C_4n^{1/6-\eta/2}
             .
         \end{split}
     \end{equation}
   
    Plugging these upper bounds into the estimate of $J_1$ obtained above, we conclude that
    \begin{align}
        &\frac{J_1}{C^*} = \quad \frac{1}{C^*} \sum_{(k,l)\in K_1\cup K_2\cup K_4\cup K_5}  |\mathcal{L}^{0}_n(X_{k,l})| e^{-\frac{X_{k,l}(a,b)}{2}}\\
        & \leq (n^{t/2-1/6+\delta-\eta}+1)  n^{-\frac{1}{6}-(\eta-\delta)}
        + n^{1/3+\delta-\eta} n^{(1-p-t)/2}
        + n^{1/6-\eta/2} n^{-1/3} + \ n^{1/3+\delta-\eta} n^{-1/3}\\
        & \leq n^{t/2-1/3}+n^{-(p+t)/2+5/6+(\delta-\eta)} + n^{-1/6-\eta/2} + n^{\delta-\eta},
    \end{align}
    where  $C^*$ is positive and does not depend on $a, \, b, \, n$. Recall that by choice of $p,t$ (see \eqref{eq:p_t_cond}), we have  $p+t>\tfrac{5}{3}-2(\eta-\delta)$ and $t<\tfrac{2}{3}$. Hence, each of the exponents of $n$ in the last estimate is negative. This implies that
    \begin{equation}
        \lim\limits_{n\to\infty} \sup\limits_{(a,b) \in S_n(\delta, \eta)}\sum_{(k,l)\in K_1\cup K_2\cup K_4\cup K_5}
        |\mathcal{L}^{0}_n(X_{k,l}(a,b))| e^{-\frac{X_{k,l}(a,b)}{2}} = 0.
    \end{equation}

    It remains to prove that the same holds for the sum over $K_3$. Let us rewrite this region as
    \begin{equation}
        K_3 = \{(k,l)\in K: n^{1/3} \leq s - X_{k,l} < n^{p}\}.
    \end{equation}
    We use a dyadic decomposition: define
    \begin{equation}
        M=\lceil \log_2 n^{p-1/3}\rceil , \qquad U_j:=2^j n^{1/3}, \,\,\ j=0,\dots,M,
    \end{equation}
    and split $K_3$ into 
    \begin{equation}
        K_3:=\bigcup\limits_{j=0}^{M-1} V_j := \bigcup\limits_{j=0}^{M-1} \left\{ (k,l) \in K_3: s-X_{k,l} \in [U_j, U_{j+1})  \right\}.
    \end{equation}
    For each layer $j$ and $(k,l)\in V_j$, for sufficiently large $n$, we have
    \begin{equation}
        X_{k,l}-q = s-q-(s-X_{k,l}(a,b)) \geq s-q-n^p \geq n+2, \quad s-X_{k,l}\geq U_j.
    \end{equation}
    Together with the first estimate in~\eqref{eq:estimate}, this gives us
    \begin{equation}\label{eq:Laguerre_on_dyadic_layer}
        |\mathcal{L}_n^{0}(X_{k,l}(a,b))|e^{-\frac{X_{k,l}(a,b)}{2}} \leq  \sqrt{\frac{s-q}{(s-X_{k,l})(X_{k,l}-q)}} \leq \frac{2}{\sqrt{U_j}}.
    \end{equation}

    Next, we estimate the number of lattice points in each layer $V_j$. We use Lemma~\ref{lem:ellipse_lattice_count} and inequality~\eqref{eq:ellipse_difference}:
    \begin{equation}
        \#V_j \leq \#((E_{s-U_j}\setminus E_{s-U_{j+1}})\cap \Z^2) \leq ab (U_{j+1}-U_{j}) + \widetilde{C} \left(a \sqrt{s-U_j} + b \sqrt{U_{j+1}-U_{j}} +1\right),
    \end{equation}
    where $j=0,\dots,M-1$ and $\widetilde{C}$ does not depend on $a, \, b, \, n, \, j$. Since $U_{j+1}-U_j = U_j = 2^j n^{1/3}$ we obtain, using~\eqref{eq:ab_bounds_revisited},
    \begin{equation}
        \#V_j \leq 2^j n^{-1/3-\eta} + \widetilde{C} (n^{1/6-\eta/2} + 2^{j/2} n^{\delta-\eta} +1).
    \end{equation}
    Combining this with~\eqref{eq:Laguerre_on_dyadic_layer}, we get
    \begin{align}
        \sum\limits_{(k,l)\in K_3} |\mathcal{L}_n^{0}(X_{k,l}(a,b))|e^{-\frac{X_{k,l}(a,b)}{2}}
        & = \sum\limits_{j=0}^{M-1} \sum\limits_{(k,l)\in V_j}
        |\mathcal{L}_n^{0}(X_{k,l}(a,b))|e^{-\frac{X_{k,l}(a,b)}{2}}\\
        & \leq \widetilde{C} \sum\limits_{j=0}^{M-1} \frac{\#V_j}{U_j^{1/2}}\\
        & \leq \widetilde{C} \sum\limits_{j=0}^{M-1}  \frac{2^j n^{-1/3-\eta} + n^{1/6-\eta/2} + 2^{j/2} n^{\delta-\eta}}{2^{j/2}n^{1/6}}. 
    \end{align}
    We treat the latter sums separately:
    \begin{equation}
        \Sigma_1 := \sum\limits_{j=0}^{M-1} 2^{j/2}n^{-1/2-\eta}
        \leq \frac{2^{M/2}}{\sqrt{2}-1} n^{-1/2-\eta}
        \leq 4n^{p/2-2/3-\eta},
    \end{equation}
    \begin{equation}
        \Sigma_2 := \sum\limits_{j=0}^{M-1} 2^{-j/2}n^{-\eta/2}
        \leq 4n^{-\eta/2},
    \end{equation}
    \begin{equation}
        \Sigma_3 := \sum\limits_{j=0}^{M-1} n^{-1/6+\delta-\eta}
        \leq M n^{-1/6+\delta-\eta}
         \leq C n^{-1/6+\delta-\eta} \log_2 n.
    \end{equation}
    Since $p<1$ and  $0 < \delta < \eta$, we see that
    \begin{equation}
        \lim_{n \to \infty} \Sigma_\ell = 0, \quad \ell = 1,2,3.
    \end{equation}
Thus, 
\begin{equation}
     \lim\limits_{n \to \infty} \sup\limits_{(a,b) \in S_n(\delta, \eta)} \sum\limits_{(k,l)\in K} |\mathcal{L}_n^{0}(X_{k,l}(a,b))|e^{-\frac{X_{k,l}(a,b)}{2}} \leq \lim_{n \to \infty} \sup\limits_{(a,b) \in S_n(\delta, \eta)} (J_1 + \Sigma_1 +\Sigma_2+ \Sigma_3) = 0.
\end{equation}
This finishes the proof of Proposition~\ref{pro:Janssen_main}.
\end{proof}

\end{document}
